% Generated draft: requires mathematical and source review.
\documentclass[11pt]{article}
\usepackage[margin=25mm]{geometry}
\usepackage{fontspec}
\setmainfont{Latin Modern Roman}
\usepackage{amsmath,amssymb,mathtools,unicode-math}
\setmathfont{Latin Modern Math}
\usepackage{enumitem,needspace,xurl,hyperref}
\setlength{\parindent}{0pt}
\setlength{\parskip}{5pt}
\setlength{\emergencystretch}{3em}
\newcommand{\N}{\mathbb{N}}
\newcommand{\Z}{\mathbb{Z}}
\newcommand{\Q}{\mathbb{Q}}
\newcommand{\R}{\mathbb{R}}
\newcommand{\C}{\mathbb{C}}
\newcommand{\F}{\mathbb{F}}
\newcommand{\A}{\mathbb{A}}
\newcommand{\PP}{\mathbb{P}}
\newcommand{\E}{\mathbb{E}}
\newcommand{\eps}{\varepsilon}
\newcommand{\dd}{\,\mathrm{d}}
\begin{document}
\section*{UnsolvedMath ID: 800268}
\textbf{Source problem code:} catalogue entry
\subsection{Definitions and mathematical statement}
Let \(\PP^2(\C)\) denote the complex projective plane, whose points are complex lines in \(\C^3\). Interpret the supplied phrase ``a holomorphic map on \(\mathbb{P}^2\) of degree \(d\)'' as a holomorphic self-map
\[
f\colon \PP^2(\C)\longrightarrow\PP^2(\C)
\]
of algebraic degree \(d\). Under this convention, \(f\) has a representation
\[
f([z_0:z_1:z_2])=[F_0(z):F_1(z):F_2(z)],
\]
where the \(F_i\) are homogeneous complex polynomials of common degree \(d\) with no common zero in \(\C^3\setminus\{0\}\). The complex ground field, self-map interpretation, and algebraic degree convention are implicit in the supplied wording rather than explicitly specified there.

A point \(p\in\PP^2(\C)\) is a repelling fixed point of \(f\) if \(f(p)=p\) and every eigenvalue \(\lambda\) of the complex derivative
\[
D f_p\colon T_p\PP^2(\C)\longrightarrow T_p\PP^2(\C)
\]
satisfies \(|\lambda|>1\).

The supplied question, with these conventions, is: for every integer \(d\geq2\) and every holomorphic self-map \(f\) of algebraic degree \(d\), must there exist \(p\in\PP^2(\C)\) such that
\[
f(p)=p
\quad\text{and}\quad
\operatorname{Spec}(D f_p)\subset\{\lambda\in\C:|\lambda|>1\}?
\]
\subsection{Short English statement}
Does every holomorphic self-map of the complex projective plane of algebraic degree at least two have a fixed point at which both derivative eigenvalues have modulus strictly greater than one? The complex self-map and degree conventions make explicit the interpretation used in this draft.
\subsection{Sources}
\begin{itemize}
\item Buzzard, Fornaess, Krantz, \emph{The Several Complex Variables Problem List} (1995), Problem 67. Supplied source URL: \url{https://arxiv.org/abs/math/9509202}. The source record reports access on 2026-10-16 and extraction by \texttt{arxiv\_latex}; these provenance statements are supplied metadata, not independently verified access.
\item Local UnsolvedMath registry, permanent ID \(800268\), source problem code \texttt{catalogue entry}, title ``Several Complex Variables Problem List --- Problem 67,'' in the AMR Open Problem Lists collection. The registry attributes the problem to Buzzard, Fornaess, Krantz, with proposed year 1995, and labels it \texttt{open}; this is a supplied status label, not an independently verified current mathematical assessment.
\end{itemize}
\end{document}
